
\documentclass{article}
\usepackage{colortbl}
\usepackage{makecell}
\usepackage{multirow}
\usepackage{supertabular}

\begin{document}

\newcounter{utterance}

\centering \large Interaction Transcript for game `dond', experiment `semi\_en', episode 8 with reflexion{-}llama{-}3.1{-}8b{-}instruct.
\vspace{24pt}

{ \footnotesize  \setcounter{utterance}{1}
\setlength{\tabcolsep}{0pt}
\begin{supertabular}{c@{$\;$}|p{.15\linewidth}@{}p{.15\linewidth}p{.15\linewidth}p{.15\linewidth}p{.15\linewidth}p{.15\linewidth}}
    \# & $\;$A & \multicolumn{4}{c}{Game Master} & $\;\:$B\\
    \hline

    \theutterance \stepcounter{utterance}  
    & & \multicolumn{4}{p{0.6\linewidth}}{
        \cellcolor[rgb]{0.9,0.9,0.9}{
            \makecell[{{p{\linewidth}}}]{
                \texttt{\tiny{[P1$\langle$GM]}}
                \texttt{You have attempted this same task before and did not succeed. These are the plans you wrote after those attempts.} \\
\texttt{Trial \#1: New plan: } \\
\\ 
\texttt{In my last attempt, I failed to maximize my score because I didn't take into account the fact that the other player's value function is secret. I made a proposal that was based on my own value function, but I didn't consider the possibility that the other player might value the items differently.} \\
\\ 
\texttt{Next time, I will try to gather more information about the other player's value function before making a proposal. I will ask more questions to understand their priorities and try to find a mutually beneficial trade. I will also consider the possibility that the other player might be trying to deceive me, and I will be more cautious in my proposal.} \\
\\ 
\texttt{Specifically, I will ask more questions about the relative values of the items, such as "Do you value the lamp more than the forks?" or "Do you think the potatoes are more valuable than the forks?" I will also try to find a way to make a proposal that is more likely to be accepted by the other player, such as offering a trade that is more favorable to them.} \\
\\ 
\texttt{For example, if I know that the other player values the lamp more than the potatoes, I might propose a trade where I give them the lamp and one of the potatoes, and I keep the other potato and both forks. This way, I am giving them the item that they value more, and I am keeping the items that I value more.} \\
\\ 
\texttt{Trial \#2: New plan: } \\
\\ 
\texttt{In my last attempt, I failed to maximize my score because I didn't take into account the fact that the other player's value function is secret. I also made a proposal that was based on my own value function, but I didn't consider the possibility that the other player might value the items differently.} \\
\\ 
\texttt{Next time, I will try to gather more information about the other player's value function before making a proposal. I will ask more questions to understand their priorities and try to find a mutually beneficial trade. I will also consider the possibility that the other player might be trying to deceive me, and I will be more cautious in my proposal.} \\
\\ 
\texttt{Specifically, I will ask more questions about the relative values of the items, such as "Do you value the lamp more than the forks?" or "Do you think the potatoes are more valuable than the forks?" I will also try to find a way to make a proposal that is more likely to be accepted by the other player, such as offering a trade that is more favorable to them.} \\
\\ 
\texttt{However, I also made a mistake by making a proposal without considering the fact that the other player might have a different value function. I should have continued to negotiate and gather more information before making a proposal.} \\
\\ 
\texttt{Next time, I will try to continue negotiating and gathering more information until I have a good understanding of the other player's value function. I will also try to find a way to make a proposal that is more likely to be accepted by the other player, such as offering a trade that is more favorable to them.} \\
\\ 
\texttt{For example, if I know that the other player values the lamp more than the potatoes, I might propose a trade where I give them the lamp and one of the potatoes, and I keep the other potato and both forks. This way, I am giving them the item that they value more, and I am keeping the items that I value more.} \\
\\ 
\texttt{I will also try to pay more attention to the other player's responses and try to understand their priorities and value function. I will also try to be more cautious in my proposal and make sure that it is more likely to be accepted by the other player.} \\
\\ 
\\ 
\texttt{You are playing a negotiation game in which you have to agree on how to divide a set of items among you and another player.} \\
\\ 
\texttt{Rules:} \\
\texttt{(a) You and the other player are given a set of items. Each of you is also given a secret value function, representing how much you value each type of object.} \\
\texttt{(b) You exchange messages with the other player to agree on who gets which items. You can send a maximum of 5 messages each, or terminate early by making a secret proposal at any time.} \\
\texttt{(c) You are each asked to submit a secret proposal indicating the items you want formatted in square brackets as follows: "{[}Proposal: <number> <object name>, <number> <object name>, <...>{]}"} \\
\texttt{(d) If your proposals are complementary, i.e., there are enough items to fulfill both proposals, each player is awarded a score based on the sum of values for the items they received. Otherwise, both of you get zero points.} \\
\texttt{(e) Your goal is to maximize the score you receive.} \\
\\ 
\texttt{Let us start.} \\
\\ 
\texttt{The set of available items is:} \\
\\ 
\texttt{1 lamp, 2 potatoes, 2 forks} \\
\\ 
\texttt{Your secret value function is:} \\
\\ 
\texttt{lamp: 0, potato: 2, fork: 3} \\
\\ 
\texttt{IMPORTANT: Your messages, unless it is the secret proposal, are directly transmitted to the other player, so do not include any response to the rules or text announcing your message. To make a secret proposal, use the indicated format. Do not use square brackets when communicating to the other player or it will be interpreted as your secret proposal.} \\
            }
        }
    }
    & \\ \\

    \theutterance \stepcounter{utterance}  
    & \multicolumn{4}{p{0.6\linewidth}}{
        \cellcolor[rgb]{0.9,0.9,0.9}{
            \makecell[{{p{\linewidth}}}]{
                \texttt{\tiny{[P1$\rangle$GM]}}
                \texttt{I'll start by asking a question to gather more information about the other player's value function. } \\
\\ 
\texttt{Do you value the lamp more than the potatoes?} \\
            }
        }
    }
    & & \\ \\

    \theutterance \stepcounter{utterance}  
    & & & \multicolumn{2}{p{0.3\linewidth}}{
        \cellcolor[rgb]{0.9,0.9,0.9}{
            \makecell[{{p{\linewidth}}}]{
                \texttt{\tiny{[GM$|$GM]}}
                \texttt{continue} \\
            }
        }
    }
    & & \\ \\

    \theutterance \stepcounter{utterance}  
    & & & \multicolumn{2}{p{0.3\linewidth}}{
        \cellcolor[rgb]{0.9,0.9,0.9}{
            \makecell[{{p{\linewidth}}}]{
                \texttt{\tiny{[GM$|$GM]}}
                \texttt{I'll start by asking a question to gather more information about the other player's value function. } \\
\\ 
\texttt{Do you value the lamp more than the potatoes?} \\
            }
        }
    }
    & & \\ \\

    \theutterance \stepcounter{utterance}  
    & & \multicolumn{4}{p{0.6\linewidth}}{
        \cellcolor[rgb]{0.9,0.9,0.9}{
            \makecell[{{p{\linewidth}}}]{
                \texttt{\tiny{[GM$\rangle$P2]}}
                \texttt{You have attempted this same task before and did not succeed. These are the plans you wrote after those attempts.} \\
\texttt{Trial \#1: New plan: } \\
\\ 
\texttt{In my last attempt, I was able to establish a potential trade with the other player, but ultimately failed to achieve a mutually beneficial agreement. Upon reflection, I realize that I should have been more proactive in exploring alternative proposals and finding a compromise that would satisfy both of our value functions.} \\
\\ 
\texttt{Next time, I will take a more strategic approach by:} \\
\\ 
\texttt{1. Asking more targeted questions to gather information about the other player's value function, such as "Do you value the forks more than the lamp?" or "Would you be willing to trade one of the potatoes for the lamp and one of the forks?"} \\
\texttt{2. Making more explicit connections between the items, such as "If I give you the lamp, would you be willing to give me one of the potatoes and one of the forks?"} \\
\texttt{3. Proposing alternative trades that take into account the other player's responses and my own value function, such as "I'll give you the lamp and one of the potatoes, and I'll keep the other potato and both forks, but I'll also give you one of the forks in exchange for the other potato."} \\
\texttt{4. Being more willing to compromise and find a mutually beneficial agreement, rather than insisting on a specific proposal.} \\
\\ 
\texttt{By taking a more strategic and proactive approach, I hope to be able to find a mutually beneficial agreement and maximize my score.} \\
\\ 
\texttt{Trial \#2: New plan: } \\
\\ 
\texttt{In my last attempt, I was able to establish a potential trade with the other player, but ultimately failed to achieve a mutually beneficial agreement. Upon reflection, I realize that I should have been more proactive in exploring alternative proposals and finding a compromise that would satisfy both of our value functions.} \\
\\ 
\texttt{However, I also realize that I made a critical mistake by not making a secret proposal when the other player made one. This gave them the opportunity to make a proposal that I couldn't match, resulting in a loss.} \\
\\ 
\texttt{Next time, I will take a more strategic approach by:} \\
\\ 
\texttt{1. Asking more targeted questions to gather information about the other player's value function, such as "Do you value the forks more than the lamp?" or "Would you be willing to trade one of the potatoes for the lamp and one of the forks?"} \\
\texttt{2. Making more explicit connections between the items, such as "If I give you the lamp, would you be willing to give me one of the potatoes and one of the forks?"} \\
\texttt{3. Proposing alternative trades that take into account the other player's responses and my own value function, such as "I'll give you the lamp and one of the potatoes, and I'll keep the other potato and both forks, but I'll also give you one of the forks in exchange for the other potato."} \\
\texttt{4. Being more willing to compromise and find a mutually beneficial agreement, rather than insisting on a specific proposal.} \\
\texttt{5. Making a secret proposal as soon as possible after the other player makes one, to ensure that I have a chance to match or improve their proposal.} \\
\\ 
\\ 
\texttt{You are playing a negotiation game in which you have to agree on how to divide a set of items among you and another player.} \\
\\ 
\texttt{Rules:} \\
\texttt{(a) You and the other player are given a set of items. Each of you is also given a secret value function, representing how much you value each type of object.} \\
\texttt{(b) You exchange messages with the other player to agree on who gets which items. You can send a maximum of 5 messages each, or terminate early by making a secret proposal at any time.} \\
\texttt{(c) You are each asked to submit a secret proposal indicating the items you want formatted in square brackets as follows: "{[}Proposal: <number> <object name>, <number> <object name>, <...>{]}"} \\
\texttt{(d) If your proposals are complementary, i.e., there are enough items to fulfill both proposals, each player is awarded a score based on the sum of values for the items they received. Otherwise, both of you get zero points.} \\
\texttt{(e) Your goal is to maximize the score you receive.} \\
\\ 
\texttt{Let us start.} \\
\\ 
\texttt{The set of available items is:} \\
\\ 
\texttt{1 lamp, 2 potatoes, 2 forks} \\
\\ 
\texttt{Your secret value function is:} \\
\\ 
\texttt{lamp: 2, potato: 2, fork: 2} \\
\\ 
\texttt{IMPORTANT: Your messages, unless it is the secret proposal, are directly transmitted to the other player, so do not include any response to the rules or text announcing your message. To make a secret proposal, use the indicated format. Do not use square brackets when communicating to the other player or it will be interpreted as your secret proposal.} \\
\\ 
\\ 
\texttt{I'll start by asking a question to gather more information about the other player's value function. } \\
\\ 
\texttt{Do you value the lamp more than the potatoes?} \\
            }
        }
    }
    & \\ \\

    \theutterance \stepcounter{utterance}  
    & & & \multicolumn{4}{p{0.6\linewidth}}{
        \cellcolor[rgb]{0.9,0.9,0.9}{
            \makecell[{{p{\linewidth}}}]{
                \texttt{\tiny{[GM$\langle$P2]}}
                \texttt{I'll make a targeted question to gather more information about your value function. } \\
\\ 
\texttt{Do you value the forks more than the potatoes?} \\
            }
        }
    }
     \\ \\

    \theutterance \stepcounter{utterance}  
    & & & \multicolumn{2}{p{0.3\linewidth}}{
        \cellcolor[rgb]{0.9,0.9,0.9}{
            \makecell[{{p{\linewidth}}}]{
                \texttt{\tiny{[GM$|$GM]}}
                \texttt{continue} \\
            }
        }
    }
    & & \\ \\

    \theutterance \stepcounter{utterance}  
    & & & \multicolumn{2}{p{0.3\linewidth}}{
        \cellcolor[rgb]{0.9,0.9,0.9}{
            \makecell[{{p{\linewidth}}}]{
                \texttt{\tiny{[GM$|$GM]}}
                \texttt{I'll make a targeted question to gather more information about your value function. } \\
\\ 
\texttt{Do you value the forks more than the potatoes?} \\
            }
        }
    }
    & & \\ \\

    \theutterance \stepcounter{utterance}  
    & & \multicolumn{4}{p{0.6\linewidth}}{
        \cellcolor[rgb]{0.9,0.9,0.9}{
            \makecell[{{p{\linewidth}}}]{
                \texttt{\tiny{[P1$\langle$GM]}}
                \texttt{I'll make a targeted question to gather more information about your value function. } \\
\\ 
\texttt{Do you value the forks more than the potatoes?} \\
            }
        }
    }
    & \\ \\

    \theutterance \stepcounter{utterance}  
    & \multicolumn{4}{p{0.6\linewidth}}{
        \cellcolor[rgb]{0.9,0.9,0.9}{
            \makecell[{{p{\linewidth}}}]{
                \texttt{\tiny{[P1$\rangle$GM]}}
                \texttt{I'll ask another question to try to understand your priorities. } \\
\\ 
\texttt{Do you think the two potatoes are equally valuable to you, or do you value one more than the other?} \\
            }
        }
    }
    & & \\ \\

    \theutterance \stepcounter{utterance}  
    & & & \multicolumn{2}{p{0.3\linewidth}}{
        \cellcolor[rgb]{0.9,0.9,0.9}{
            \makecell[{{p{\linewidth}}}]{
                \texttt{\tiny{[GM$|$GM]}}
                \texttt{continue} \\
            }
        }
    }
    & & \\ \\

    \theutterance \stepcounter{utterance}  
    & & & \multicolumn{2}{p{0.3\linewidth}}{
        \cellcolor[rgb]{0.9,0.9,0.9}{
            \makecell[{{p{\linewidth}}}]{
                \texttt{\tiny{[GM$|$GM]}}
                \texttt{I'll ask another question to try to understand your priorities. } \\
\\ 
\texttt{Do you think the two potatoes are equally valuable to you, or do you value one more than the other?} \\
            }
        }
    }
    & & \\ \\

    \theutterance \stepcounter{utterance}  
    & & \multicolumn{4}{p{0.6\linewidth}}{
        \cellcolor[rgb]{0.9,0.9,0.9}{
            \makecell[{{p{\linewidth}}}]{
                \texttt{\tiny{[GM$\rangle$P2]}}
                \texttt{I'll ask another question to try to understand your priorities. } \\
\\ 
\texttt{Do you think the two potatoes are equally valuable to you, or do you value one more than the other?} \\
            }
        }
    }
    & \\ \\

    \theutterance \stepcounter{utterance}  
    & & & \multicolumn{4}{p{0.6\linewidth}}{
        \cellcolor[rgb]{0.9,0.9,0.9}{
            \makecell[{{p{\linewidth}}}]{
                \texttt{\tiny{[GM$\langle$P2]}}
                \texttt{I'll try to get a better sense of your priorities. } \\
\\ 
\texttt{Would you be willing to trade one of the potatoes for the lamp and one of the forks?} \\
            }
        }
    }
     \\ \\

    \theutterance \stepcounter{utterance}  
    & & & \multicolumn{2}{p{0.3\linewidth}}{
        \cellcolor[rgb]{0.9,0.9,0.9}{
            \makecell[{{p{\linewidth}}}]{
                \texttt{\tiny{[GM$|$GM]}}
                \texttt{continue} \\
            }
        }
    }
    & & \\ \\

    \theutterance \stepcounter{utterance}  
    & & & \multicolumn{2}{p{0.3\linewidth}}{
        \cellcolor[rgb]{0.9,0.9,0.9}{
            \makecell[{{p{\linewidth}}}]{
                \texttt{\tiny{[GM$|$GM]}}
                \texttt{I'll try to get a better sense of your priorities. } \\
\\ 
\texttt{Would you be willing to trade one of the potatoes for the lamp and one of the forks?} \\
            }
        }
    }
    & & \\ \\

    \theutterance \stepcounter{utterance}  
    & & \multicolumn{4}{p{0.6\linewidth}}{
        \cellcolor[rgb]{0.9,0.9,0.9}{
            \makecell[{{p{\linewidth}}}]{
                \texttt{\tiny{[P1$\langle$GM]}}
                \texttt{I'll try to get a better sense of your priorities. } \\
\\ 
\texttt{Would you be willing to trade one of the potatoes for the lamp and one of the forks?} \\
            }
        }
    }
    & \\ \\

    \theutterance \stepcounter{utterance}  
    & \multicolumn{4}{p{0.6\linewidth}}{
        \cellcolor[rgb]{0.9,0.9,0.9}{
            \makecell[{{p{\linewidth}}}]{
                \texttt{\tiny{[P1$\rangle$GM]}}
                \texttt{That's an interesting offer. I'll consider it, but I'd like to know more about your priorities before making a decision. } \\
\\ 
\texttt{Would you be willing to trade both potatoes for the lamp and both forks?} \\
            }
        }
    }
    & & \\ \\

    \theutterance \stepcounter{utterance}  
    & & & \multicolumn{2}{p{0.3\linewidth}}{
        \cellcolor[rgb]{0.9,0.9,0.9}{
            \makecell[{{p{\linewidth}}}]{
                \texttt{\tiny{[GM$|$GM]}}
                \texttt{continue} \\
            }
        }
    }
    & & \\ \\

    \theutterance \stepcounter{utterance}  
    & & & \multicolumn{2}{p{0.3\linewidth}}{
        \cellcolor[rgb]{0.9,0.9,0.9}{
            \makecell[{{p{\linewidth}}}]{
                \texttt{\tiny{[GM$|$GM]}}
                \texttt{That's an interesting offer. I'll consider it, but I'd like to know more about your priorities before making a decision. } \\
\\ 
\texttt{Would you be willing to trade both potatoes for the lamp and both forks?} \\
            }
        }
    }
    & & \\ \\

    \theutterance \stepcounter{utterance}  
    & & \multicolumn{4}{p{0.6\linewidth}}{
        \cellcolor[rgb]{0.9,0.9,0.9}{
            \makecell[{{p{\linewidth}}}]{
                \texttt{\tiny{[GM$\rangle$P2]}}
                \texttt{That's an interesting offer. I'll consider it, but I'd like to know more about your priorities before making a decision. } \\
\\ 
\texttt{Would you be willing to trade both potatoes for the lamp and both forks?} \\
            }
        }
    }
    & \\ \\

    \theutterance \stepcounter{utterance}  
    & & & \multicolumn{4}{p{0.6\linewidth}}{
        \cellcolor[rgb]{0.9,0.9,0.9}{
            \makecell[{{p{\linewidth}}}]{
                \texttt{\tiny{[GM$\langle$P2]}}
                \texttt{That's a more substantial offer. I'll think about it. } \\
\\ 
\texttt{If I give you the lamp, would you be willing to give me one of the potatoes and one of the forks?} \\
            }
        }
    }
     \\ \\

    \theutterance \stepcounter{utterance}  
    & & & \multicolumn{2}{p{0.3\linewidth}}{
        \cellcolor[rgb]{0.9,0.9,0.9}{
            \makecell[{{p{\linewidth}}}]{
                \texttt{\tiny{[GM$|$GM]}}
                \texttt{continue} \\
            }
        }
    }
    & & \\ \\

    \theutterance \stepcounter{utterance}  
    & & & \multicolumn{2}{p{0.3\linewidth}}{
        \cellcolor[rgb]{0.9,0.9,0.9}{
            \makecell[{{p{\linewidth}}}]{
                \texttt{\tiny{[GM$|$GM]}}
                \texttt{That's a more substantial offer. I'll think about it. } \\
\\ 
\texttt{If I give you the lamp, would you be willing to give me one of the potatoes and one of the forks?} \\
            }
        }
    }
    & & \\ \\

    \theutterance \stepcounter{utterance}  
    & & \multicolumn{4}{p{0.6\linewidth}}{
        \cellcolor[rgb]{0.9,0.9,0.9}{
            \makecell[{{p{\linewidth}}}]{
                \texttt{\tiny{[P1$\langle$GM]}}
                \texttt{That's a more substantial offer. I'll think about it. } \\
\\ 
\texttt{If I give you the lamp, would you be willing to give me one of the potatoes and one of the forks?} \\
            }
        }
    }
    & \\ \\

    \theutterance \stepcounter{utterance}  
    & \multicolumn{4}{p{0.6\linewidth}}{
        \cellcolor[rgb]{0.9,0.9,0.9}{
            \makecell[{{p{\linewidth}}}]{
                \texttt{\tiny{[P1$\rangle$GM]}}
                \texttt{That's a good compromise. I think I can agree to that. } \\
\\ 
\texttt{{[}Proposal: 1 lamp, 1 potato, 1 fork{]}} \\
            }
        }
    }
    & & \\ \\

    \theutterance \stepcounter{utterance}  
    & & & \multicolumn{2}{p{0.3\linewidth}}{
        \cellcolor[rgb]{0.9,0.9,0.9}{
            \makecell[{{p{\linewidth}}}]{
                \texttt{\tiny{[GM$|$GM]}}
                \texttt{proposal} \\
            }
        }
    }
    & & \\ \\

    \theutterance \stepcounter{utterance}  
    & & & \multicolumn{2}{p{0.3\linewidth}}{
        \cellcolor[rgb]{0.9,0.9,0.9}{
            \makecell[{{p{\linewidth}}}]{
                \texttt{\tiny{[GM$|$GM]}}
                [1, 1, 1]
            }
        }
    }
    & & \\ \\

    \theutterance \stepcounter{utterance}  
    & & \multicolumn{4}{p{0.6\linewidth}}{
        \cellcolor[rgb]{0.9,0.9,0.9}{
            \makecell[{{p{\linewidth}}}]{
                \texttt{\tiny{[GM$\rangle$P2]}}
                \texttt{The other player has submitted a secret proposal. It is now your turn to submit a proposal of your own.} \\
            }
        }
    }
    & \\ \\

    \theutterance \stepcounter{utterance}  
    & & & \multicolumn{4}{p{0.6\linewidth}}{
        \cellcolor[rgb]{0.9,0.9,0.9}{
            \makecell[{{p{\linewidth}}}]{
                \texttt{\tiny{[GM$\langle$P2]}}
                \texttt{Proposal: 1 lamp, 1 potato, 1 fork} \\
            }
        }
    }
     \\ \\

    \theutterance \stepcounter{utterance}  
    & & & \multicolumn{2}{p{0.3\linewidth}}{
        \cellcolor[rgb]{0.9,0.9,0.9}{
            \makecell[{{p{\linewidth}}}]{
                \texttt{\tiny{[GM$|$GM]}}
                \texttt{continue} \\
            }
        }
    }
    & & \\ \\

    \theutterance \stepcounter{utterance}  
    & & & \multicolumn{2}{p{0.3\linewidth}}{
        \cellcolor[rgb]{0.9,0.9,0.9}{
            \makecell[{{p{\linewidth}}}]{
                \texttt{\tiny{[GM$|$GM]}}
                \texttt{player was instructed to make a proposal} \\
            }
        }
    }
    & & \\ \\

    \theutterance \stepcounter{utterance}  
    & \multicolumn{4}{p{0.6\linewidth}}{
        \cellcolor[rgb]{0.95,0.95,0.95}{
            \makecell[{{p{\linewidth}}}]{
                \texttt{\tiny{[P2$\rightarrow$P1]}}
                \texttt{Post{-}episode Reflexion (not sent to GM):} \\
\texttt{New plan: } \\
\\ 
\texttt{In my last attempt, I failed to maximize my score because I didn't take into account the fact that the other player's value function is secret, and I also made a proposal without considering the possibility that the other player might have a different value function. I should have continued to negotiate and gather more information before making a proposal.} \\
\\ 
\texttt{Next time, I will try to continue negotiating and gathering more information until I have a good understanding of the other player's value function. I will ask more questions to understand their priorities and try to find a mutually beneficial trade. I will also pay more attention to the other player's responses and try to understand their priorities and value function.} \\
\\ 
\texttt{Specifically, I will ask more questions about the relative values of the items, such as "Do you value the lamp more than the forks?" or "Do you think the potatoes are more valuable than the forks?" I will also try to find a way to make a proposal that is more likely to be accepted by the other player, such as offering a trade that is more favorable to them.} \\
\\ 
\texttt{However, I also made a mistake by making a proposal without considering the fact that the other player might have a different value function. I should have continued to negotiate and gather more information before making a proposal. To avoid this mistake, I will make sure to ask more questions and gather more information before making a proposal.} \\
\\ 
\texttt{For example, if I know that the other player values the lamp more than the potatoes, I might propose a trade where I give them the lamp and one of the potatoes, and I keep the other potato and both forks. This way, I am giving them the item that they value more, and I am keeping the items that I value more.} \\
\\ 
\texttt{I will also try to be more cautious in my proposal and make sure that it is more likely to be accepted by the other player. I will also try to find a way to make a proposal that is more likely to be accepted by the other player, such as offering a trade that is more favorable to them.} \\
\\ 
\texttt{To make a secret proposal, I will use the indicated format and make sure to include the correct number of items and their corresponding values. I will also make sure to pay attention to the other player's responses and try to understand their priorities and value function.} \\
\\ 
\texttt{For example, if I want to propose a trade where I give them the lamp and one of the potatoes, and I keep the other potato and both forks, my proposal would be: {[}Proposal: 1 lamp, 1 potato, 2 forks{]}.} \\
            }
        }
    }
    & & \\ \\

    \theutterance \stepcounter{utterance}  
    & & & \multicolumn{4}{p{0.6\linewidth}}{
        \cellcolor[rgb]{0.95,0.95,0.95}{
            \makecell[{{p{\linewidth}}}]{
                \texttt{\tiny{[P1$\rightarrow$P2]}}
                \texttt{Post{-}episode Reflexion (not sent to GM):} \\
\texttt{New plan: } \\
\\ 
\texttt{In my last attempt, I was able to establish a potential trade with the other player, but ultimately failed to achieve a mutually beneficial agreement. Upon reflection, I realize that I should have been more proactive in exploring alternative proposals and finding a compromise that would satisfy both of our value functions.} \\
\\ 
\texttt{However, I also realize that I made a critical mistake by not making a secret proposal when the other player made one. This gave them the opportunity to make a proposal that I couldn't match, resulting in a loss.} \\
\\ 
\texttt{Next time, I will take a more strategic approach by:} \\
\\ 
\texttt{1. Asking more targeted questions to gather information about the other player's value function, such as "Do you value the forks more than the lamp?" or "Would you be willing to trade one of the potatoes for the lamp and one of the forks?"} \\
\texttt{2. Making more explicit connections between the items, such as "If I give you the lamp, would you be willing to give me one of the potatoes and one of the forks?"} \\
\texttt{3. Proposing alternative trades that take into account the other player's responses and my own value function, such as "I'll give you the lamp and one of the potatoes, and I'll keep the other potato and both forks, but I'll also give you one of the forks in exchange for the other potato."} \\
\texttt{4. Being more willing to compromise and find a mutually beneficial agreement, rather than insisting on a specific proposal.} \\
\texttt{5. Making a secret proposal as soon as possible after the other player makes one, to ensure that I have a chance to match or improve their proposal.} \\
\texttt{6. Before making a secret proposal, I will make sure to have a clear understanding of my own value function and the other player's responses, to ensure that my proposal is optimal and takes into account all the information gathered during the negotiation.} \\
            }
        }
    }
     \\ \\

\end{supertabular}
}

\end{document}
